\subsection{Related Work}
\label{sec:related_work}
\providecommand{\finalchange}[1]{#1}
\providecommand{\localexecchange}[1]{#1}

Language-guided grounding methods translate natural-language instructions into
spatial representations, executable programs, or scene-grounded task
plans~\cite{huang2023inner,liang2023code,rana2023sayplan,
huang2023voxposer}. Building on these capabilities, LLM-based multi-robot
frameworks adopt different coordination structures: RoCo uses inter-robot
dialogue; COHERENT relies on centralized proposals and execution feedback;
MALLVI reactivates specialized agents when needed; and CoMuRoS and DEXTER-LLM
revise tasks or assignments in response to events~\cite{
mandi2024roco,liu2025coherent,taji2026mallvi,
borate2026comuros,zhu2025dexter}. Other systems explicitly incorporate
capability-aware coalitions, dependency-based task allocation, or separate
planning, manipulation, and verification agents~\cite{
kannan2024smartllm,obata2025lipllm,he2026closedloop}. Collectively, these
methods establish effective mechanisms for grounding, agent interaction, and
task revision, but they do not explicitly organize ownership of the distinct
semantic components that must remain consistent as multiple heterogeneous
manipulators execute concurrently. RoSC addresses this issue by maintaining
the authoritative shared state and assigning responsibility
to the affordance, planning, and verification
roles, respectively. When a discrete semantic event occurs, these roles
identify its dependency-closed affected scope and reconstruct only the
invalidated subgraph, while preserving unaffected valid execution and leaving
continuous physical adaptation to the local controllers.


Dynamic manipulation supports tracked or reactive grasping and visual
servoing~\cite{marturi2019dynamic,maranci2026reactive,liu2023target,
haviland2020finalphase}, allowing local controllers to continuously adapt
end-effector motion to moving targets and perception updates. Execution
monitoring and continual replanning further address runtime deviations by
detecting discrepancies between expected and observed outcomes and invoking
recovery or plan revision~\cite{pettersson2005execution,brenner2009continual,
levihn2013foresight,pezzato2023active}. These methods primarily focus on
maintaining physical feasibility, monitoring action execution, or revising a
plan after deviations. However, in heterogeneous multi-manipulator settings,
the resulting feedback may affect different parts of the team-level semantic
state: a new observation may change an object's affordance, a failed action may
invalidate its execution status and downstream dependencies, and a revised
instruction may alter only one task branch. RoSC explicitly separates these
responsibilities among affordance, verification, and planning roles. It uses
the shared semantic state to determine which information must be refreshed,
which subtasks are affected, and whether planning should be invoked, thereby
supporting event-scoped semantic revision while preserving unaffected valid
execution.
